\documentclass[conference,10pt]{IEEEtran}

\usepackage{amsmath,amssymb}
\usepackage{graphicx}
\usepackage{booktabs}
\usepackage{cite}
\usepackage{url}
\usepackage[hidelinks]{hyperref}

\begin{document}

\title{LyASC: Stability-Driven Skip-Connection Gating\\for Deep Residual Networks}

\author{\IEEEauthorblockN{Anonymous Authors}
\IEEEauthorblockA{Department of Computer Science\\
Email: authors@example.com}}

\maketitle

% ===================================================================
%  ABSTRACT
% ===================================================================
\begin{abstract}
Deep residual networks usually fix the weights of their skip
connections at design time. Scaling those weights away from unity hurts
accuracy unless a normalization layer follows the combination, and even
the best fixed scaling schedules have to be tuned by hand for each
dataset. A freely learned gate would remove the tuning, but it
underperforms the fixed schedules because its parameters receive no
structural signal during training. We propose LyASC (Lyapunov-Aware
Skip Connections), which drives the gate with an online, per-block
estimate of local dynamical instability. Power iteration approximates
the dominant singular value of each residual block's Jacobian, a cheap
proxy for the finite-time Lyapunov exponent, and a sigmoid turns that
value into a per-block coefficient that blends the skip and residual
signals. This is a partial report covering CIFAR-10 with
PreAct-ResNet-110. Under normal training, LyASC reaches 7.17\% test
error. That beats the unconstrained learned gate (8.16\%) but not the
best fixed baseline (6.15\%). Under stress conditions that push fixed
baselines to 37\% to 49\% error, LyASC brings training back below 9\%
error in two of the three regimes we tested.
\end{abstract}


% ===================================================================
%  I. INTRODUCTION
% ===================================================================
\section{Introduction}

Deep residual networks~\cite{he2016deep} pass information between
layers through skip connections, which add each block's input directly
to its output. He et al.~\cite{he2016identity} showed that this identity
mapping is enough to train networks with more than 100 layers. Later
work has asked whether other skip weightings can do better.

Liu et al.~\cite{liu2021rethinking} evaluated skip-connection scaling on
PreAct-ResNets and reported two findings that matter here. First,
scaling the skip path by a fixed factor $\lambda > 1$ (that is,
$y = \lambda x + F(x)$) hurts performance when $\lambda$ is large, but
adding Layer Normalization after the combination
($y = \text{LN}(\lambda x + F(x))$) recovers baseline accuracy and
sometimes exceeds it. Their best configuration, a recursively normalized
variant (2rSkip+LN), reached 6.02\% error on CIFAR-10 with
PreAct-ResNet-110. Second, replacing the fixed $\lambda$ with a learned
per-channel weight vector (wSkip+LN) did worse, at 8.66\% error, and the
learned weights shrank to small magnitudes during training.

This leaves a gap. Fixed schedules work well, but they need tuning for
each dataset and cannot adapt to the local training dynamics of
individual blocks. The unconstrained learned gate fails because it gets
no signal about when a given block would benefit from favoring the skip
path over the residual branch.

LyASC (Lyapunov-Aware Skip Connections) replaces both with a gate
driven by a measured quantity: the local dynamical instability of each
residual block. At every block, a power-iteration estimate of the
dominant singular value of the residual branch's Jacobian serves as a
proxy for the finite-time Lyapunov exponent. A sigmoid gate turns this
estimate into a blending coefficient $\alpha_l \in (0, 1)$ that sets the
balance between the skip and residual paths. When the residual branch is
locally unstable (large singular value), the gate moves toward the skip
path. When the branch is stable, the gate lets more of the residual
signal through.

The contributions of this paper are:
\begin{enumerate}
    \item The LyASC architecture: stability-driven skip gating that
    fits into standard PreAct-ResNet blocks without changing the
    backbone interface.
    \item A verified reproduction of the base paper's image
    classification results on CIFAR-10. Of 32 rows, 14 pass a
    $\pm$0.4 percentage-point tolerance, and every non-catastrophic
    deviation is an outperform.
    \item Normal-condition results. LyASC with post-gate Layer
    Normalization reaches 7.17\% error. It beats the unconstrained
    learned gate (8.16\%) by 0.99 percentage points and falls 1.02
    percentage points short of the best fixed baseline (6.15\%).
    \item Stress-test results. Under skip scaling that drives fixed
    baselines to 37\% to 49\% error, LyASC brings every seed below 9\%
    error in two of three regimes. The third regime gives a partially
    resolved two-cluster outcome.
\end{enumerate}

This report covers the implementation, stress-testing, and normalization pairing ablations (Phase~E2) of LyASC. CIFAR-100 depth scaling and systematic hyperparameter searches remain as future work.


% ===================================================================
%  II. RELATED WORK
% ===================================================================
\section{Related Work}

\subsection{Residual Learning, Skip Scaling, and Normalization}

He et al.~\cite{he2016deep, he2016identity} established identity skip
connections as the standard mechanism enabling stable optimization of deep
residual networks exceeding 100 layers. Zaeemzadeh et al.~\cite{zaeemzadeh2020norm}
analyzed norm-preservation properties of residual connections, showing that
preserving signal and gradient norms across depth prevents gradient vanishing.
Liu et al.~\cite{liu2021rethinking} systematically evaluated fixed skip-scaling
parameters alongside post-combination Layer Normalization in ResNets and
Transformers, discovering that fixed recursive scaling (2rSkip+LN) achieves
competitive performance (6.02\% on CIFAR-10), whereas unconstrained learned
scaling (wSkip+LN) degrades to 8.66\%. MacDonald et al.~\cite{macdonald2023skip}
further clarified the interaction between skip paths and normalization layers in
deep optimization landscapes, showing that normalization rescales activations
by a static geometric rule without tracking dynamical branch behavior during
training. Veit et al.~\cite{veit2016residual} conceptualized residual networks as
ensembles of relatively shallow paths, implying that modulating skip and
residual contributions alters the effective depth distribution of the ensemble.

\subsection{Learned Gating and Dynamic Routing}

Highway Networks~\cite{srivastava2015highway} introduced learned transform and
carry gates ($T(x)$ and $C(x) = 1 - T(x)$) to dynamically modulate feature flow
across layers. Subsequent work explored dynamic execution paths: SkipNet~\cite{wang2018skipnet}
and adaptive inference graph architectures~\cite{veit2020adaptive} formulate
sample-dependent routing policies primarily to reduce inference compute.
Huang et al.~\cite{huang2016deep} introduced stochastic depth, randomly bypassing
residual blocks during training according to a predefined linear decay schedule.
More recently, Joseph et al.~\cite{joseph2024lambda} proposed Lambda-Skip connections
to prevent rank collapse in deep representations, while Thota~\cite{thota2025cosinegate}
introduced CosineGate for semantic dynamic routing based on directional cosine
incompatibility. Unlike content-based gates that condition on feature activations
or stochastic schedules, LyASC is driven by an online dynamical stability signal
computed from local Jacobian amplification.

\subsection{Dynamical Systems Perspectives and Edge of Stability}

A growing body of work examines deep network training through dynamical systems
and stability theory. Nar and Sastry~\cite{nar2018residual} established Lyapunov
stability conditions and convex decomposition for residual networks, treating
layer-to-layer activations as discrete-time dynamical trajectories.
Cohen et al.~\cite{cohen2021gradient} demonstrated that gradient descent on
neural networks typically operates at the ``edge of stability'' (EoS), where
the top eigenvalue of the training loss Hessian hovers around $2/\eta$.
Subsequent theoretical analyses by Arora et al.~\cite{arora2022understanding},
Damian et al.~\cite{damian2022self}, and Andreyev and Beneventano~\cite{andreyev2024edge}
further characterized implicit self-stabilization and stochastic stability under SGD.
Mittra~\cite{mittra2024utilizing} and Sharafi et al.~\cite{sharafi2024weight}
explored Lyapunov exponents and weight trajectory metrics for architecture design
and learning dynamics diagnostics.
Ba et al.~\cite{ba2016layer} demonstrated that Layer Normalization stabilizes
hidden-state dynamics by standardizing activation variance.
LyASC bridges these domains by employing an online, finite-difference Jacobian-vector
product (JVP) amplification proxy as a local controller signal to dynamically
modulate the skip-residual blending parameter $\alpha_l$.

\subsection{Literature Comparative Analysis}

Table~\ref{tab:lit_comparison} summarizes how established research families
relate to and differ from LyASC.

\begin{table}[t]
\centering
\caption{Research families and their structural relationship to LyASC.}
\label{tab:lit_comparison}
\begin{tabular}{p{2.2cm}p{2.6cm}p{2.8cm}}
\toprule
\textbf{Research Family} & \textbf{Established Contribution} & \textbf{Distinction from LyASC} \\
\midrule
Residual identity mappings~\cite{he2016deep, he2016identity} & Direct signal and gradient route across depth & Fixed identity contribution without adaptive modulation \\
\addlinespace
Skip scaling \& norm~\cite{liu2021rethinking, macdonald2023skip} & Demonstrates scaling sensitivity and norm interaction & Static hand-tuned schedules rather than online proxy \\
\addlinespace
Norm-preservation~\cite{zaeemzadeh2020norm} & Analyzes depth trainability via isometric bounds & Theoretical framework; does not gate active branches \\
\addlinespace
Learned gating~\cite{srivastava2015highway, hu2018squeeze} & Learns parametric branch contributions & Feature-driven rather than stability-signal-driven \\
\addlinespace
Dynamic routing~\cite{wang2018skipnet, veit2020adaptive} & Selects execution paths by input or policy & Focuses on compute efficiency rather than gradient stability \\
\addlinespace
Stochastic depth~\cite{huang2016deep} & Random block dropping during training & Stochastic drop schedule; not signal-guided \\
\addlinespace
EoS analysis~\cite{cohen2021gradient, arora2022understanding, damian2022self} & Characterizes global Hessian optimization dynamics & Passive diagnostic; does not actively control layer gates \\
\addlinespace
Lyapunov analysis~\cite{nar2018residual, mittra2024utilizing} & Quantifies trajectory perturbation growth & LyASC deploys a finite local proxy as real-time gate controller \\
\bottomrule
\end{tabular}
\end{table}



% ===================================================================
%  III. PROPOSED METHOD
% ===================================================================
\section{Proposed Method}

\subsection{Generalized Residual-Block Formulation}

In a standard PreAct-ResNet block $l$ with input $x_l$ and residual
function $F_l(x_l; \theta_l)$, the output is $y_l = x_l + F_l(x_l)$. Liu
et al.~\cite{liu2021rethinking} generalize this to
$y_l = G(\lambda x_l + F_l(x_l))$, where $\lambda$ is a fixed scalar
and $G$ is an optional normalization function.

LyASC formalizes a generalized stability-modulated residual block:
\begin{align}
    \alpha_l &= \sigma\bigl(\beta_l \cdot (\hat{\lambda}_l - \lambda_0)\bigr), \label{eq:gate} \\
    y_l &= G\bigl(\alpha_l \cdot s_x x_l + (1 - \alpha_l) \cdot s_F F_l(x_l; \theta_l)\bigr) \label{eq:output}
\end{align}
where $s_x$ and $s_F$ denote configured branch scale factors, $G$ is the
post-gate normalization operator, and $\alpha_l \in (0, 1)$ acts as a dynamic
convex blending parameter when $s_x = s_F = 1$. Under standard unstressed
training, $s_x = s_F = 1$. In stress testing, $s_x$ and $s_F$ are varied to
reproduce structural scaling collapse modes (e.g., $s_x=0.5$ or $s_F \in \{2, 3\}$).

Here $\hat{\lambda}_l$ is an online exponentially smoothed local stability
proxy at block~$l$ (Section~\ref{sec:estimator}), $\lambda_0 = 0$ is the
nominal stability threshold, $\beta_l$ is a per-block learnable sensitivity
parameter initialized to 0, and $\sigma(z) = 1 / (1 + e^{-z})$ is the standard
logistic sigmoid. Because $\beta_l = 0$ at initialization, $\alpha_l = 0.5$ for
every block at step zero, ensuring that every block starts with an equal nominal
blend of identity and residual pathways before sensitivity parameters diverge.
The parameter $\beta_l$ is not constrained to be strictly positive; the network
learns the directional sensitivity per block via backpropagation.

We investigate three post-gate operations: $G \in \{\text{Identity}, \text{LayerNorm2d}, \text{BatchNorm2d}\}$.
Under primary evaluation, $G = \text{LayerNorm2d}$ is adopted following the
ablation in Section~\ref{sec:ablation1}.

\subsection{Stability Estimator as Local Amplification Proxy}
\label{sec:estimator}

The controller signal $\hat{\lambda}_l$ approximates the local amplification of
perturbations across the residual branch $F_l$. Let $J_l(x_l) = \partial F_l / \partial x_l$
denote the local Jacobian of block $l$. A singular value $\sigma_{\max}(J_l) > 1$
indicates local perturbation expansion, whereas $\sigma_{\max}(J_l) < 1$ indicates
contraction.

A full finite-time Lyapunov exponent (FTLE) calculation requires tracking perturbation
growth across an extended temporal trajectory of operators with continuous orthogonalization.
In our implementation, we compute an online, finite-difference Jacobian-vector product
(JVP) amplification proxy using $k$-step power iteration without computing the full Jacobian:
\begin{enumerate}
    \item Sample a random unit Gaussian probe vector $v \in \mathbb{R}^d$, normalized such that $\|v\|_2 = 1$.
    \item Compute the directional directional amplification $u = J_l \cdot v$ via a single JVP in PyTorch autograd.
    \item Update the probe vector direction $v \leftarrow u / (\|u\|_2 + \epsilon)$.
    \item Repeat steps 2 and 3 for $k$ power-iteration steps (configured to $k=1$ in our primary pipeline).
    \item Define the raw logarithmic amplification proxy:
    \begin{equation}
        \lambda_l^{\text{raw}} = \log\bigl(\|u\|_2 + \epsilon\bigr).
    \end{equation}
\end{enumerate}

To dampen batch-level gradient noise and high-frequency stochasticity, the raw proxy is
smoothed via an exponential moving average (EMA):
\begin{equation}
    \hat{\lambda}_l^{(t)} = \rho \cdot \hat{\lambda}_l^{(t-1)} + (1 - \rho) \cdot \lambda_l^{\text{raw},(t)}
\end{equation}
where $\rho = 0.9$ is the EMA decay factor. To bound computational overhead, the power-iteration
probe is updated once every $N$ training iterations ($N = 20$), and the cached smoothed
statistic $\hat{\lambda}_l$ is read during intermediate forward passes.

\subsection{Controller Interpretation and System Scope}

From a dynamical systems and control perspective, LyASC operates as a feedback-inspired
local controller: it samples local amplification, converts it to a convex mixing gate,
and adjusts branch throughput dynamically as optimization progresses. We emphasize that
LyASC provides an empirical heuristic rather than a formal Lyapunov stability guarantee:
it does not compute global state-space Lyapunov spectra, nor does it guarantee bounded
asymptotic dynamics under arbitrary nonlinear perturbations. Instead, it offers a
computationally tractable proxy ($k=1, N=20$) that alerts the architecture when a
residual branch is severely amplifying feature magnitudes.

\subsection{Integration with PreAct-ResNet-110}

LyASC is implemented via two self-contained modules in \texttt{lyasc\_model}:
\begin{itemize}
    \item \texttt{StabilityEstimator}: Manages the persistent probe buffer, EMA smoothing, and $N$-step cadence.
    \item \texttt{LyASCSkip}: Implements Eqs.~\ref{eq:gate} and \ref{eq:output}, exposing the standard \texttt{forward(x, F\_x)} interface expected by \texttt{PreActBasicBlock}.
\end{itemize}
The baseline codebase (\texttt{base\_model}) is maintained as a strictly read-only reference.
The sensitivity parameter $\beta_l$ is excluded from weight decay, matching the treatment
accorded to normalization affine scales and the $w$ parameter of wSkip+LN.


% ===================================================================
%  IV. EXPERIMENTAL SETUP
% ===================================================================
\section{Experimental Setup}

All experiments use CIFAR-10~\cite{krizhevsky2009learning}, which has
50,000 training images in 10 classes and 10,000 test images. Training
augmentation follows the standard protocol: zero-pad 4 pixels per side,
take a random $32 \times 32$ crop, flip horizontally at random, and
normalize per channel. We apply no augmentation at test time.

The backbone is PreAct-ResNet-110~\cite{he2016identity}, which has 54
two-layer residual blocks in three stages with channel widths 16, 32,
and 64. Downsampling at stage transitions uses Option~A (zero-padded
identity shortcut with average-pool stride), as in the base-model
reproduction.

Training uses SGD with momentum~0.9, weight decay $2 \times 10^{-4}$,
and batch size~128. The learning rate starts at 0.1 and is divided by 10
at iterations 32,000 and 48,000, and training ends at 64,000 iterations
($\approx$128 epochs). We clip gradients
($\texttt{max\_norm} = 5.0$) at every step, a choice carried over from
the base-model reproduction, where LN-augmented variants needed it to
train stably. The affine parameters of Batch
Normalization~\cite{ioffe2015batch} and Layer
Normalization~\cite{ba2016layer}, all biases, and the per-block
$\beta_l$ gate sensitivity parameters are excluded from weight decay.

Each configuration is evaluated over three random seeds. The base paper
used five. We used fewer because of compute limits and list this as a
limitation. All comparisons use the same symmetric $\pm 0.4$
percentage-point tolerance as the base-model reproduction.

LyASC's estimator hyperparameters are $k = 1$ (power-iteration steps
per recomputation), $N = 20$ (training steps between recomputations),
EMA decay $\rho = 0.9$, and stability threshold $\lambda_0 = 0.0$. We
apply no cold-start handling: the estimator runs from iteration~0 with
no warmup or damping (Section~\ref{sec:phased} reports the ablation
behind this choice).

\subsection{Hardware Environment and Compute Constraints}

All training runs execute on a dedicated single-GPU workstation with a
4GB VRAM limit. A full 64,000-iteration training run using LayerNorm2d
allocates up to 3.8~GB of VRAM, which prevents overlapping execution and
requires all training runs to be strictly serialized. Training uses
standard single-precision floating-point arithmetic (FP32) without
mixed precision.

\subsection{Software Environment and Automated Verification}

The experimental pipeline is isolated in an environment running Python~3.10
and PyTorch with CUDA acceleration. The base-model codebase
(\texttt{base\_model/}) is maintained as a strictly read-only reference.
To verify implementation correctness, an automated test suite comprising
78 unit tests across \texttt{lyasc\_model/tests/} and \texttt{base\_model/tests/}
checks Jacobian-vector product calculations, exponential moving average updates,
buffer persistence across training iterations, checkpoint restoration, and
gradient propagation without modifying the baseline modules.

\subsection{Experiment Tracking and Data Provenance}

All experimental runs, hyperparameters, training trajectories (cross-entropy
loss, classification accuracy, gradient norms, and per-block gate statistics),
and checkpoint locations are logged to an MLflow tracking service backed by
a local SQLite database. All numerical results reported in this work trace
directly to saved result JSON artifacts and serialized checkpoints.



% ===================================================================
%  V. RESULTS
% ===================================================================
\section{Results}

\subsection{Base-Model Reproduction}
\label{sec:reproduction}

Before evaluating LyASC, we reproduced the CIFAR-10 results of
Liu et al.~\cite{liu2021rethinking} (Tables~1, 3, 4, and~6 on
PreAct-ResNet-110). Of 32 rows, 14 pass the $\pm 0.4$
percentage-point tolerance. Every non-catastrophic deviation is an
outperform, meaning our error rate is strictly lower than the paper's.
The three catastrophic Table~3 rows (LN($0.5x + F$), LN($x + 2F$),
LN($x + 3F$)) reproduce the instability and the bimodal per-seed
outcomes the paper describes. Table~\ref{tab:baselines} lists the
reference values used in the rest of this paper.

\begin{table}[t]
\centering
\caption{Selected base-model reproduction results on CIFAR-10,
PreAct-ResNet-110 (3 seeds). ``Paper'': value reported by
Liu et al.~\cite{liu2021rethinking}. Tolerance: $\pm$0.4~pp.}
\label{tab:baselines}
\begin{tabular}{lccl}
\toprule
Variant & Ours (\%) & Paper (\%) & Verdict \\
\midrule
1xSkip        & 6.27 $\pm$ 0.04 & 6.31  & PASS \\
2rSkip+LN     & 6.15 $\pm$ 0.09 & 6.02  & PASS \\
wSkip+LN      & 8.16 $\pm$ 0.13 & 8.66  & FAIL$^*$ \\
\midrule
LN($0.5x+F$) & 37.14 $\pm$ 9.94  & 29.88 & FAIL$^\dagger$ \\
LN($x+2F$)   & 37.59 $\pm$ 22.83 & 24.98 & FAIL$^\dagger$ \\
LN($x+3F$)   & 49.27 $\pm$ 15.31 & 61.02 & FAIL$^\dagger$ \\
\bottomrule
\multicolumn{4}{l}{\footnotesize $^*$Outperform (our number is lower).} \\
\multicolumn{4}{l}{\footnotesize $^\dagger$Catastrophic regime;
instability mechanism reproduced.}
\end{tabular}
\end{table}

We found and fixed two implementation bugs during the reproduction. The
first was a LayerNorm axis error: the code normalized across the full
spatial extent instead of channel-wise per position. The second was a
weight-decay scope error: the \texttt{wSkip+LN} gate parameter
\texttt{w} was being decayed because its 3-D shape slipped past the
standard dimensionality exclusion check. Both are recorded in the
project's decisions log.


\subsection{Phase C: Initial LyASC Evaluation}
\label{sec:phasec}

The initial LyASC configuration (no post-gate normalization,
$G = \text{identity}$) reached a 3-seed mean error of 9.66\% $\pm$
0.74 (individual seeds: 10.21\%, 8.82\%, 9.95\%). That is a higher error
than both the best fixed baseline (2rSkip+LN at 6.15\%) and the
unconstrained learned gate (wSkip+LN at 8.16\%).

The likely cause is the missing normalization after the gate. Without
it, the convex combination $\alpha_l x + (1 - \alpha_l) F(x)$ has a
magnitude that depends on $\alpha_l$, and that shift compounds over 54
consecutive blocks. The final gate distribution was wide ($\alpha$
standard deviation of 0.296 across blocks at convergence), which made
the effect worse. This diagnosis led to the LayerNorm ablation below.


\subsection{LayerNorm After the Gate (Ablation 1)}
\label{sec:ablation1}

Adding $G = \text{LayerNorm2d}$ (channel-wise Layer Normalization) after
the gate output cut test error from 10.21\% to 7.22\% on seed~0, a drop
of 2.99 percentage points. This was a single-seed proof of concept, and
the 3-seed evaluation in Phase~D (Section~\ref{sec:phased}) confirmed
the benefit.

With LayerNorm, the gate distribution tightened. The $\alpha$ standard
deviation fell from 0.296 to 0.201, and mean $\alpha$ rose from 0.598 to
0.681, so blocks converged toward more similar gate values instead of
spreading across the full $[0, 1]$ range. The 7.22\% error beats
wSkip+LN (8.16\%) by 0.94 percentage points and cuts the gap to
2rSkip+LN (6.15\%) from 4.06 to 1.07 percentage points.

We adopted LyASC+LN as the default configuration for all later
experiments. The empirical benefit of $G = \text{LayerNorm2d}$ is clear.
How it improves gate behavior (by controlling activation scale,
preventing saturation, or through some other effect) was not measured
and remains open.


\subsection{Cold-Start Strategy Comparison (Phase D)}
\label{sec:phased}

In the base-model reproduction, LN-augmented skip variants showed large
gradient norm spikes (hundreds to low thousands) during the first 1,000
iterations, and the stable variants recovered by iteration~2,000.
LyASC's estimator sees the same transient. We compared three strategies
for handling it:

\begin{itemize}
    \item \textbf{None:} the estimator is active from iteration~0 with
    no special handling.
    \item \textbf{Damped:} the estimator is active, but its influence
    on~$\alpha$ ramps linearly from 0 to 1 over the first 1,000
    iterations.
    \item \textbf{Warmup:} $\alpha$ is held at 0.5 for the first 1,000
    iterations, and the estimator output is ignored during that window.
\end{itemize}

\begin{table}[t]
\centering
\caption{Cold-start strategy comparison (LyASC+LN,
PreAct-ResNet-110, CIFAR-10). $n$: number of seeds.}
\label{tab:coldstart}
\begin{tabular}{lccc}
\toprule
Strategy & $n$ & Mean error (\%) & Std (\%) \\
\midrule
None     & 3 & 7.17 & 0.21 \\
Damped   & 3 & 7.46 & 0.31 \\
Warmup   & 2 & 7.95 & 0.21 \\
\bottomrule
\end{tabular}
\end{table}

Table~\ref{tab:coldstart} shows the results. At $n = 3$, ``none'' and
``damped'' cannot be told apart: the 0.29 percentage-point gap between
their means is inside the project's $\pm$0.4~pp tolerance band and
smaller than the standard deviation of ``damped'' (0.31~pp). ``Warmup''
is worst at 7.95\% (2 seeds only), which suggests that switching the
estimator off completely during early training hurts.

We used ``none'' in later experiments because it needs no extra
hyperparameter (ramp duration or warmup length) and its mean error is at
least as low as the alternatives. The per-seed values for ``none'' were
7.33\%, 7.25\%, and 6.93\%.


\subsection{Stress-Test Evaluation (Phase E1)}
\label{sec:stress}

The stress tests scale the skip path or the residual path by factors
that cause catastrophic training failure in the fixed baselines
reproduced from Table~3 of Liu et al.~\cite{liu2021rethinking}, and
check whether LyASC's gate can compensate. Every stress configuration
uses the standard LyASC+LN setup, the ``none'' cold-start strategy, and
the hyperparameters used in all other experiments.

\begin{table}[t]
\centering
\caption{Stress-test results (LyASC+LN, sample SD). ``Baseline'': the
corresponding fixed-weight Table~3 result from the base-model
reproduction. $n$: number of seeds.}
\label{tab:stress}
\begin{tabular}{lccl}
\toprule
Condition & $n$ & LyASC+LN (\%) & Baseline (\%) \\
\midrule
stress\_0.5 & 3 & 8.38 $\pm$ 0.43 & 37.14 \\
stress\_2.0 & 3 & 8.44 $\pm$ 0.34 & 37.59 \\
stress\_3.0 & 5 & see text         & 49.27 \\
\bottomrule
\end{tabular}
\end{table}

Table~\ref{tab:stress} summarizes the results. In stress\_0.5 (skip path
scaled by 0.5, matching the LN($0.5x+F$) regime) and stress\_2.0
(residual branch scaled by~2, matching LN($x+2F$)), every seed converges
to below~9\% error, while the fixed baselines fail at 37\% to 38\%. The
per-seed values are 7.98\%, 8.84\%, and 8.31\% for stress\_0.5, and
8.13\%, 8.81\%, and 8.38\% for stress\_2.0.

The stress\_3.0 condition (residual branch scaled by~3, matching
LN($x+3F$)) produced two clusters across 5~seeds:

\begin{itemize}
    \item \textbf{Full rescue} (1 of 5 seeds): seed~4 reached 8.75\%
    error with a final mean gate value of $\alpha = 0.83$.
    \item \textbf{Partial rescue} (4 of 5 seeds): mean error 25.16\%,
    median 20.88\%, range 19.12\% to 32.73\%, with final $\alpha$ values
    near 0.64 to 0.66.
\end{itemize}

All five seeds beat the fixed baseline (49.27\%), but the outcomes
differ widely from seed to seed. A pooled 5-seed mean of 21.88\% $\pm$
9.15 would mislead, because it averages two qualitatively different
regimes.

The per-iteration gate trajectories suggest a hypothesis: how fast a
seed escapes in the first $\sim$16,000 iterations decides which cluster
it ends up in. Seed~4's gate passed $\alpha = 0.71$ by iteration~16,000
and kept rising to 0.83. The four partial-rescue seeds moved more
slowly, and their $\alpha$ plateaued near 0.65 and never rose further.
Within the partial-rescue cluster, final error tracks when the plateau
was escaped relative to the learning-rate drop at iteration~32,000.
Seeds that escaped by iteration~27,000 to 28,000 finished at 19\% to
21\% error, and seeds that stalled or relapsed at the LR transition
finished at 28\% to 33\%.

This threshold-effect account fits the trajectory data, but it is a
hypothesis and not a confirmed mechanism. The data do not settle whether
slow gate progression caused incomplete convergence, or whether the
3$\times$ scaling made convergence hard on its own and the gate merely
tracked the difficulty. Estimating how often full and partial rescue
occur at this stress level requires more seeds.


\subsection{Normalization Pairing Ablations (Phase E2)}
\label{sec:phasee2}

To evaluate how the choice of post-gate normalization function $G$ affects performance, Phase~E2 tested two normalization variants on PreAct-ResNet-110 (2 seeds each, 64,000 iterations), comparing them against the closed-out LyASC+LN Phase~D baseline (7.17\% $\pm$ 0.21\%):
\begin{itemize}
    \item \textbf{LyASC\_noG ($G = \text{identity}$, reused from Phase~C):} Omitting post-gate normalization yields a 3-seed mean error of 9.66\% $\pm$ 0.74\% (individual seeds: 10.21\%, 8.82\%, 9.95\%), incurring a +2.49 percentage-point accuracy penalty compared to LyASC+LN (7.17\%). This confirms that post-gate normalization is necessary to maintain single-digit classification error.
    \item \textbf{LyASC+BN ($G = \text{BatchNorm2d}$):} Replacing LayerNorm with post-gate 2-D Batch Normalization yields an error of 6.34\% $\pm$ 0.01\% (individual seeds: 6.35\%, 6.33\%, $n=2$, with seed~2 in training). In unstressed training, this represents an observed difference of -0.83 percentage points compared to LyASC+LN (7.17\%, $n=3$). This 6.34\% error rate is statistically indistinguishable from plain 1xSkip (6.27\%--6.36\%) and remains above the best fixed baseline, 2rSkip+LN (6.15\%). Furthermore, the behavior of LyASC+BN under stress conditions remains untested.
\end{itemize}

Table~\ref{tab:normpairing} summarizes the normalization pairing ablation matrix.

\begin{table}[t]
\centering
\caption{Normalization pairing ablation matrix (PreAct-ResNet-110, CIFAR-10, sample SD). $n$: number of seeds.}
\label{tab:normpairing}
\begin{tabular}{lcccc}
\toprule
Variant & Post-Gate Norm $G$ & $n$ & Mean Error (\%) & $\alpha_{\text{mean}}$ \\
\midrule
LyASC+LN (Phase D)   & LayerNorm2d & 3 & 7.17 $\pm$ 0.21 & 0.701 \\
LyASC\_noG (Phase C) & Identity    & 3 & 9.66 $\pm$ 0.74 & 0.667 \\
LyASC+BN             & BatchNorm2d & 2$^*$ & 6.34 $\pm$ 0.01 & 0.562 \\
PreAct-ResNet-110    & N/A         & 3 & 6.27 $\pm$ 0.04 & N/A   \\
\bottomrule
\multicolumn{5}{l}{\footnotesize $^*$Seed 2 in training at time of reporting.}
\end{tabular}
\end{table}


% ===================================================================
%  VI. DISCUSSION AND LIMITATIONS
% ===================================================================
\section{Discussion and Limitations}

The results show that a stability-driven gate can work as a
skip-connection mechanism, with several clear limitations.

\textit{Normal-condition gap.}
LyASC+LN (7.17\%) does not beat the best fixed baseline, 2rSkip+LN
(6.15\%), under normal training. The 1.02 percentage-point gap is
outside the $\pm$0.4~pp tolerance used in this work. LyASC does beat the
unconstrained learned gate (wSkip+LN, 8.16\%), but with no stress
applied the best fixed schedule is still stronger. One reading is that
LyASC pays for generality: its gate has to find a good operating point
during training, while the fixed schedule has one built in. Whether
better hyperparameters or architectural changes can reduce that cost is
open.

\textit{Wall-clock overhead and normalization cost.}
Measured training runtimes demonstrate that computational overhead is dominated by the normalization operator rather than the stability estimator. A complete 64,000-iteration training run requires approximately 4.7 hours (0.26 s/iter) for un-normalized LyASC (Phase C), approximately 6.2 hours (0.35 s/iter) for LyASC+BN, and approximately 8.9 hours (0.50 to 0.68 s/iter) for LyASC+LN. The power-iteration estimator with $k=1$ and $N=20$ executes one single directional Jacobian-vector product every 20 iterations, contributing an estimated 3\% to 5\% overhead. The primary computational bottleneck stems from the spatial and channel reductions within the 54 LayerNorm2d operations across the deep residual trunk.

\textit{Internal validity and reproducibility safeguards.}
Several procedural safeguards protect experimental validity. The reference baseline implementation (\texttt{base\_model/}) remains completely frozen and read-only. Random number generator seeds are fixed across PyTorch, NumPy, and Python standard libraries. Quantitative results are logged automatically to an MLflow database and exported to structured JSON records, preventing clerical transcription errors.

\textit{External validity and scale of evaluation.}
All empirical evaluations in this work are conducted on a single GPU using CIFAR-10 on PreAct-ResNet-110 with three seeds per standard configuration (five seeds for stress\_3.0). The published baseline paper evaluated five seeds per configuration. We have not evaluated whether these findings generalize to larger datasets (such as ImageNet-1k), different domain modalities (such as language or audio), alternative architectural families (such as Transformers or ConvNeXt), or deeper residual networks.

\textit{Unresolved stress\_3.0 bifurcation mechanism.}
The two-cluster outcome observed under the severe stress\_3.0 condition remains incompletely resolved. While iteration trajectories show that early gate progression above $\alpha = 0.70$ before iteration 16,000 correlates with full recovery, the causal direction remains open: it is unclear whether early gate adaptation actively drives recovery or merely reflects optimization progress driven by favorable initial gradients.

\textit{Hyperparameter exploration.}
All experiments evaluated a single fixed hyperparameter configuration ($k = 1$, $N = 20$, $\lambda_0 = 0$, $\rho = 0.9$, $\beta_l$ initialized to 0). This configuration was chosen based on initial profiling rather than systematic grid search. Whether alternative thresholds $\lambda_0$, faster update intervals $N$, or multi-step power iterations $k > 1$ can close the remaining gap to the 2rSkip+LN baseline remains an open question.



% ===================================================================
%  VII. CONCLUSION AND FUTURE WORK
% ===================================================================
\section{Conclusion and Future Work}

LyASC replaces fixed or unconstrained skip-connection gates with a
per-block blending coefficient driven by an online estimate of local
dynamical instability. On CIFAR-10 with PreAct-ResNet-110, it beats the
unconstrained learned gate (wSkip+LN) by about 1 percentage point under
normal training, and it stabilizes training in two of three catastrophic
stress regimes where the fixed baselines fail outright. It does not
match the best hand-tuned fixed baseline, 2rSkip+LN.

Three key directions remain. First, Phase~F will investigate
novel architectural extensions, including dynamic $\lambda$ learning rate
schedulers and adaptive EMA decay rates. Evaluating depth scaling on
CIFAR-100 at depths 20 through 110 remains as separate future work.
Second, a systematic search over the estimator hyperparameters ($k$,
$N$, $\lambda_0$, and the initialization of~$\beta_l$) may close the
1.02~pp gap to the fixed baseline under normal conditions. Third, the
project's research plan lists ten proposed extensions to the LyASC
formulation, including per-channel gating, hierarchical multi-scale
estimators, and explicit coupling to the training loss Hessian, that go
beyond the single-scalar, per-block gate evaluated here.


% ===================================================================
%  REFERENCES
% ===================================================================
\bibliographystyle{IEEEtran}
\bibliography{references}

\end{document}